% =============================================================================
%  导言区：中译注本《同调镜像对称与环面纤维化》
%  原文：arXiv:math/0011041v2 [math.SG], 4 Jun 2001, 66 pp.
%        Maxim Kontsevich, Yan Soibelman
%        "Homological mirror symmetry and torus fibrations"
%  编译：xelatex main.tex（三遍；两遍定目录、一遍定交叉引用）
% =============================================================================
\documentclass[11pt,a4paper]{article}

\usepackage[UTF8, fontset=fandol, heading=true]{ctex}
\usepackage{amsmath}
\usepackage{amssymb}
\usepackage{amsthm}
\usepackage{mathtools}
\usepackage{bm}
\usepackage[a4paper,left=2.5cm,right=2.5cm,top=2.6cm,bottom=2.6cm]{geometry}
\usepackage{graphicx}
\graphicspath{{figures/}}
\usepackage{xcolor}
\usepackage[most]{tcolorbox}
\usepackage{enumitem}
\usepackage{booktabs}
\usepackage{longtable}
\usepackage{array}
\usepackage{tikz}
\usepackage{tikz-cd}
\usetikzlibrary{shapes.geometric,arrows.meta,positioning,calc}
\usepackage[font=small,labelfont=bf,labelsep=quad]{caption}
\usepackage{cite}
\usepackage[colorlinks=true,
            linkcolor=blue!45!black,
            citecolor=blue!45!black,
            urlcolor=blue!45!black]{hyperref}

\linespread{1.15}
\setlength{\emergencystretch}{3em}
\allowdisplaybreaks[3]
\pagestyle{plain}
\renewcommand{\figurename}{图}
\renewcommand{\refname}{参考文献}
\renewcommand{\contentsname}{目录}
\renewcommand{\appendixname}{附录}
\renewcommand{\proofname}{证明}
\renewcommand{\abstractname}{摘要}

% 原文用 \emph 标出强调，中文以加粗对应
\renewcommand{\emph}[1]{{\bfseries #1}}

% ---------------------------------------------------------------------------
% 常用记号（与原文一致；译手只允许使用这里定义的宏与 amsmath 标准命令）
% ---------------------------------------------------------------------------
\newcommand{\A}{\mathcal{A}}
\newcommand{\B}{\mathcal{B}}
\newcommand{\D}{\mathcal{D}}
\newcommand{\E}{\mathcal{E}}
\newcommand{\F}{\mathcal{F}}
\newcommand{\G}{\mathcal{G}}
\newcommand{\HH}{\mathcal{H}}
\newcommand{\LL}{\mathcal{L}}
\newcommand{\MM}{\mathcal{M}}
\newcommand{\NN}{\mathcal{N}}
\newcommand{\OO}{\mathcal{O}}
\newcommand{\cP}{\mathcal{P}}
\newcommand{\SSS}{\mathcal{S}}
\newcommand{\TT}{\mathcal{T}}
\newcommand{\UU}{\mathcal{U}}
\newcommand{\VV}{\mathcal{V}}
\newcommand{\WW}{\mathcal{W}}
\newcommand{\XX}{\mathcal{X}}
\newcommand{\YY}{\mathcal{Y}}
\newcommand{\ZZ}{\mathcal{Z}}
\newcommand{\Z}{\mathbb{Z}}
\newcommand{\R}{\mathbb{R}}
\newcommand{\C}{\mathbb{C}}
\newcommand{\N}{\mathbb{N}}
\newcommand{\Q}{\mathbb{Q}}
\newcommand{\PP}{\mathbb{P}}
\newcommand{\kk}{\Bbbk}
\newcommand{\eps}{\varepsilon}
\newcommand{\lam}{\lambda}
\newcommand{\p}{\partial}
\newcommand{\dbar}{\bar\partial}
\newcommand{\vol}{\mathrm{vol}}
\newcommand{\dvol}{d\,\mathrm{vol}}
\newcommand{\id}{\mathrm{id}}
\newcommand{\Id}{\mathrm{Id}}
\newcommand{\op}{\mathrm{op}}
\newcommand{\coker}{\operatorname{coker}}
\newcommand{\im}{\operatorname{im}}
\newcommand{\rank}{\operatorname{rank}}
\newcommand{\codim}{\operatorname{codim}}
\newcommand{\sgn}{\operatorname{sgn}}
\newcommand{\tr}{\operatorname{tr}}
\newcommand{\Tot}{\operatorname{Tot}}
\newcommand{\Ber}{\operatorname{Ber}}
\newcommand{\Supp}{\operatorname{Supp}}
\newcommand{\Hom}{\operatorname{Hom}}
\newcommand{\Ext}{\operatorname{Ext}}
\newcommand{\End}{\operatorname{End}}
\newcommand{\Aut}{\operatorname{Aut}}
\newcommand{\Spec}{\operatorname{Spec}}
\newcommand{\Max}{\operatorname{Max}}
\newcommand{\Ob}{\operatorname{Ob}}
\newcommand{\Fun}{\operatorname{Fun}}
\newcommand{\Coh}{\operatorname{Coh}}
\newcommand{\Perf}{\operatorname{Perf}}
\newcommand{\Mod}{\operatorname{Mod}}
\newcommand{\Fuk}{\mathsf{Fuk}}
\newcommand{\DR}{\mathrm{DR}}
\newcommand{\dR}{\mathrm{dR}}
\newcommand{\GH}{\mathrm{GH}}
\newcommand{\Fq}{\mathbb{F}_q}

% A∞ 相关简写
\newcommand{\Ainf}{A_{\infty}}
\newcommand{\ainf}{A_\infty}
\newcommand{\Aop}{A^{\mathrm{op}}}
\newcommand{\sA}{\mathfrak{A}}
\newcommand{\mB}{\mathfrak{B}}

% 原文中的黑体/花体小符号
\newcommand{\Xmer}{X^{\mathrm{mer}}_{q}}
\newcommand{\Xform}{X^{\mathrm{form}}}
\newcommand{\Cmer}{C^{\mathrm{mer}}_{q}}
\newcommand{\Ceps}{C_{\eps}}
\newcommand{\Yan}{Y^{\mathrm{an}}}
\newcommand{\TY}{T_Y}
\newcommand{\TYZ}{T^{\Z}_Y}
\newcommand{\Tstar}{T^{*}}
\newcommand{\Symp}{\mathrm{Symp}}
\newcommand{\Ham}{\mathrm{Ham}}
\newcommand{\Lag}{\mathrm{Lag}}
\newcommand{\Crit}{\mathrm{Crit}}
\newcommand{\ind}{\operatorname{ind}}
\newcommand{\grad}{\operatorname{grad}}
\newcommand{\Lie}{\mathcal{L}}

% ---------------------------------------------------------------------------
% 定理类环境：与原文编排一致——每一类环境各自连续编号（定理 1、定理 2…），
% 不带节号前缀；节号另行引用（如"§6.3"）。
% ---------------------------------------------------------------------------
\newtheoremstyle{cnth}%
  {6pt}% above
  {6pt}% below
  {\normalfont}% body
  {}% indent
  {\bfseries}% head font
  {.}% punctuation
  {0.5em}% space after head
  {}% head spec

\theoremstyle{cnth}
\newtheorem{thm}{定理}
\newtheorem{lmm}{引理}
\newtheorem{prp}{命题}
\newtheorem{crl}{推论}
\newtheorem{conj}{猜想}
\newtheorem{dfn}{定义}
\newtheorem{rmk}{注}
\newtheorem{exmp}{例}
\newtheorem{ques}{问题}
\newtheorem{cnv}{约定}
\newtheorem{nota}{记号}

% 证明末尾方框
\renewcommand{\qedsymbol}{$\square$}

% ---------------------------------------------------------------------------
% 译注方框：橙色，编号；可用 \showyizhufalse 一键关闭（关闭后文字以内文排印）
% ---------------------------------------------------------------------------
\newif\ifshowyizhu
\showyizhutrue
\newcounter{yizhu}

\newenvironment{yizhu}[1][]{%
    \refstepcounter{yizhu}%
    \par\medskip
    \ifshowyizhu
        \begin{tcolorbox}[enhanced, breakable,
            colback=orange!4!white, colframe=orange!55!black,
            boxrule=0.5pt, arc=2pt, left=6pt, right=6pt, top=4pt, bottom=4pt,
            title={\small\textbf{译注 \theyizhu\if\relax\detokenize{#1}\relax\else： #1\fi}},
            fonttitle=\bfseries, coltitle=black, colbacktitle=orange!12!white,
            toptitle=1pt, bottomtitle=1pt]
    \fi
}{%
    \ifshowyizhu
        \end{tcolorbox}%
    \fi
    \par\smallskip
}

% 译者脚注（不进入译注编号，仅用于行内小注）
\newcommand{\zy}[1]{\footnote{译者注：#1}}
